\item Una muestra de $1.00\text{ mol}$ de un gas ideal monoatómico se lleva a través del ciclo que se muestra en la figura P9.37. El proceso $A \rightarrow B$ es una expansión isotérmica reversible. Calcule (a) el trabajo neto realizado por el gas, (b) la energía agregada al gas por calor, (c) la energía expulsada del gas por calor y (d) la eficiencia del ciclo. (e) Explique cómo se compara la eficiencia con la de una máquina de Carnot que funciona entre los mismos extremos de temperatura.
